% Separate post-freeze candidate; a base-field example, not a BT comparison.
\subsection{The required cyclotomic break three is attainable}

\begin{lemma}[An explicit allowed base with the required cyclotomic quotient]
\label{lem:B-RD89-base}
\status{Candidate; complete proof supplied, independent audits pending.}
Let \(k\) be an algebraically closed field of characteristic five,
put \(K_0=W(k)[1/5]\), choose a primitive fifth root \(\zeta_5\),
and set
\[
 K_1=K_0(\zeta_5),\qquad \pi_1=\zeta_5-1.
\]
Let \(t\) be a root of
\[
       (1+T)^5-\zeta_5+\pi_1^4T^2=0,
       \qquad K=K_1(t).                              \tag{B-RD89.1}
\]
Then \(K\) is complete discretely valued, has residue field \(k\),
and satisfies \(e_K=20\) and \(\mu_5\subset K\).
The extension \(C=K(\mu_{25})/K\) is cyclic of degree five,
with lower break three. The extension \(K(\mu_{125})/K\)
is cyclic of degree 25. If \(L/K\) is the tame extension
of degree 24, then \(L(\mu_{25})/L\) has break 72.
\end{lemma}
\begin{proof}
The usual cyclotomic Eisenstein polynomial
\(\Phi_5(1+X)\) shows that \(K_1/K_0\) has degree four,
with uniformizer \(\pi_1\). Expanded over \(K_1\), the
polynomial in (B-RD89.1) is
\[
 T^5+5T^4+10T^3+(10+\pi_1^4)T^2+5T-\pi_1.
                                                        \tag{B-RD89.2}
\]
It is Eisenstein at \(\pi_1\), so \([K:K_1]=5\),
\(t\) is a uniformizer, \(v_K(\pi_1)=5\), and
\(v_K(5)=20\). Completeness and the residue-field assertion
follow from this finite totally ramified extension.

Choose \(\zeta_{25}\) with \(\zeta_{25}^5=\zeta_5\),
and put \(\alpha=\zeta_{25}-1\). The element
\(x=\alpha-t\) satisfies the exact shifted polynomial
\[
 -\pi_1^4t^2+
 \sum_{j=1}^5\binom5j(1+t)^{5-j}X^j=0.               \tag{B-RD89.3}
\]
Its constant coefficient has \(K\)-valuation 22,
its coefficients in degrees one through four have
valuation 20, and its leading coefficient is one.
Thus every intermediate point lies strictly above
the chord from \((0,22)\) to \((5,0)\). All its roots,
and in particular \(x\), have valuation \(22/5\).

Since \(\mu_5\subset K\), adjoining \(\zeta_{25}\)
gives a Galois extension of degree dividing five.
The denominator in \(v_K(x)=22/5\) forces the degree
to equal five. Normalizing \(v_C\), one has
\(v_C(x)=22\). For every nonidentity
\(g\in\operatorname{Gal}(C/K)\),
\[
  gx-x=(\zeta_5^j-1)\zeta_{25}
       \quad\text{for some }j\in\{1,2,3,4\},
  \qquad v_C(gx-x)=25.                               \tag{B-RD89.4}
\]
The leading valuation 22 is prime to five, so the
exact lower-break equality of
Lemma~\ref{lem:B-RD60-break} applies to this actual
coordinate and gives \(b(C/K)=25-22=3\).
This also proves \(\mu_{25}\not\subset K\).

For the next cyclotomic layer,
\(K_1(\mu_{125})/K_1\) is cyclic of degree 25,
with unique degree-five subfield \(K_1(\mu_{25})\).
Its intersection with the degree-five extension
\(K/K_1\) is therefore \(K_1\): a nontrivial intersection
would make \(K=K_1(\mu_{25})\), contrary to what was
just proved. Hence \(K(\mu_{125})/K\) is cyclic of
degree 25. Finally the degree-24 tame extension is
disjoint from every such cyclotomic five-extension.
The tame uniformizer/Hensel calculation written in
Corollary~\ref{cor:B-RD88-cyclotomic} gives
\(b(L(\mu_{25})/L)=24\cdot3=72\).
\end{proof}

\begin{remark}[Only the base and its cyclotomic character are realized]
\label{rem:B-RD89-scope}
\status{Candidate scope; no height-two formal group or ambiguous pair constructed.}
The example meets the cyclotomic part of the necessary
profile in B88 and gives the full order-25 cyclotomic
image at level 125. It does not construct the remaining
three independent first-layer characters, the required
height-two formal groups, or a common actual torsion
comparison. In particular it is neither a lower bound
example for the unresolved least index nor an existence
proof for the full \((4,4)\) image sector. Its role is
to rule out the unsupported shortcut that break three
could never occur over an allowed base field.

Only Eisenstein and Newton-polygon arithmetic, the
already proved exact B60 lower-break equality, and
the explicit B88 tame scaling are used. No additional
external classification theorem is imported.

The retained exact checker is
\path{code/reduction_descent/verify_p5_cyclotomic_break_three_example.py}.
The completed author run passed at 01:55:21.154128 UTC
on 20 November 2026 in 0.04096559699974023 seconds.
It verifies the two symbolic polynomial identities,
all four strict Newton vertices, and the Eisenstein
coefficients of the three cyclotomic polynomials used.
The initial execution stopped at a structural SymPy
equality comparing an expanded expression with an
unexpanded one. Its exact program, standard streams
and failure record are preserved; the corrected check
tests that the expanded polynomial difference is zero.
No mathematical assertion was inferred from the failed run.
\end{remark}
